\documentclass[11pt]{article}
\usepackage[utf8]{inputenc}
\usepackage{amsmath}
\usepackage{lipsum} % to generate some filler text
\usepackage{fullpage}
\usepackage{mathrsfs}
\usepackage{svg}
\usepackage{xr}
\usepackage{amsthm,cite,natbib,url,graphicx,booktabs,lipsum,color,bm,caption,subcaption,soul}
\usepackage{pifont,tikz,paralist,multirow,amssymb}
\usepackage{xifthen}
\usepackage{enumerate}
\usepackage{titlesec}

% 自定义 Reviewer 计数器
\newcounter{reviewer}
\setcounter{reviewer}{0}
\newcounter{point}[reviewer]
\setcounter{point}{0}

% % 将 Reviewer 的计数方式改为字母（A, B, C, D, ...）
% \renewcommand{\thereviewer}{\Alph{reviewer}}

% 定义 Reviewer 的引用格式
\renewcommand{\thepoint}{Comment\,\thereviewer.\arabic{point}}

% 定义 Associate Editor 的评论环境
\newenvironment{point1}
   {\refstepcounter{point} \bigskip \noindent {\textbf{Editor Comment~}} ---\ }
   {\par }

% 定义 Reviewer 的评论环境
\newenvironment{point}
   {\refstepcounter{point} \bigskip \noindent {\textbf{Reviewer~\thepoint} } ---\ }
   {\par }

% 定义 Reviewer 的简短评论命令
\newcommand{\shortpoint}[1]{\refstepcounter{point}   \bigskip \noindent 
	{\textbf{Reviewer~Point~\thepoint} } ---~#1  }

% 定义回复环境
\newenvironment{reply}
   {\medskip \noindent \color{blue} \begin{sf}\textbf{Reply}:\  }
   {\medskip \par \end{sf}}

% 定义简短回复命令
\newcommand{\shortreply}[2][]{\medskip \noindent \begin{sf}\textbf{Reply}:\  #2
	\ifthenelse{\equal{#1}{}}{}{ \hfill \footnotesize (#1)}%
	\medskip \end{sf}}

% 允许跳过某些 Reviewer
\newcommand{\skipreviewer}{\stepcounter{reviewer}}

% 定义 Reviewer 部分
\newcommand{\reviewersection}{\stepcounter{reviewer} \bigskip \hrule
                  \section*{Reviewer \thereviewer}}

% Title
% \title{Notes on resubmitting revision made to manuscript 259887128}
\title{Authors' Responses of the manuscript 259887128}
\author{}
\date{}

\begin{document}
\maketitle

\vspace{-30pt}

The authors would like to thank the editor and reviewers for their constructive comments and suggestions, which have helped improve the quality of this work. The revised manuscript has undergone a thorough revision according to the editor’s and reviewers’ comments. Please see our responses below. For the reviewers’ convenience, we have highlighted significant changes in the revised manuscript in \textcolor{blue}{blue}.

\section*{Editor Comments}
\begin{point1}
The reviewers see considerable merit in the manuscript, but suggest major revisions. 

Reviewers #2 and #3both provided interesting and incisive comments, and addressing them will surely improve your paper. I am not an expert on causal modelling, but the concerns raised seem reasonable and insightful to me. You may feel that the points they make can be refuted. My impression is that the concerns raised are sufficiently deep to suggest that you need to clarify the contexts in which your approach is likely to be useful, and those when it is not, and also to provide guidance for readers on the choice of the 3 parameters that the reviewer notes you rely on in demonstrating the method. How would sensible values of these parameters be chosen in a real world case as opposed to a case based on synthetic data?

Reviewer #3 offers comments that can be relatively easily addressed to improve the paper.
\end{point1}

\begin{reply}
We thank the editor for the opportunity to resubmitviseanuscript. In light of the editor’s and reviewers’ constructive suggestions, we have substantially revised the manuscript to improve its clarity, methodological rigor, and overall quality. Below, we summarize the major revisions:

\begin{itemize}
  \item In response to the reviewers’ valuable comments, we carefully re-examined the proposed GCMC method and substantially rewrote Section~2 (\textit{Materials and methods}). Specifically, we clarified the theoretical foundations of GCMC within the original method description and added a new subsection that provides detailed descriptions of the benchmark experiment setup and the three spatial cross-sectional datasets used in the case studies. Additionally, we placed the setup of the noise sensitivity analysis in this new subsection to improve the organization and clarity of the methods section.
  
  \item We thoroughly revised the \textit{Results} section. In particular, we:
  \begin{itemize}
    \item Adjusted the parameter settings for GCMC by using the false nearest neighbors method \citep{kennel1992fnn} to select an appropriate embedding dimension $E$. Following the recommendation in the original GCCM paper \citep{gao2023gccm}, we set the number of nearest neighbors $k$ for GCCM as $E + 2$, while for GCMC, $k$ was determined using a greedy search strategy that selects the value maximizing the estimated causal strength.
    \item Re-ran the three original case studies and the noise sensitivity tests using the updated settings;
    \item Introduced a new benchmark experiment to demonstrate the advantages of GCMC over GCCM.
  \end{itemize}
  
  \item In response to the reviewers’ broader concerns, we re-evaluated the overall contribution of our work and accordingly revised the \textit{Title}, \textit{Abstract}, \textit{Introduction}, \textit{Discussion}, and \textit{Conclusion} sections to better emphasize the novelty and significance of the proposed approach.
\end{itemize}

We hope that this extensive revisions address the reviewers’ concerns, and we respectfully submit the revised manuscript for review.

\end{reply}

\paragraph{Note:} All line numbers mentioned in this response letter refer to the small line numbers adjacent to the text in the revised manuscript (as generated by the \LaTeX{} line-numbering package). These may not correspond to the line numbers used by the submission system, which tend to be larger and differ from those in the manuscript.

\skipreviewer

\reviewersection


\begin{point}
The main methodological problem that I see with this method is its main idea to infer causality from consistency of spatial patterns between spatial variables. The original idea, temporal delay embedding (Takens, 1981) leverages the fact that time introduces a directed dependence—past states determine future evolution—so if variable x can reconstruct the manifold of y through time delays, it reflects causality encoded in temporal successions. Spatial embeddings, in contrast, lack such directionality: spatial coordinates capture coexistence rather than succession. consistency between spatial embeddings of x and y however only indicates similarity between spatial autocorrelation, not causal influence. Without an explicit process model describing how effects propagate across space, spatial embedding consistency captures only a necessary, but not a sufficient condition of causality. What is lacking is the control of counterfactual conditions, which only allows to infer causality (Pearl, Rubin, Peters & Schölkopf, etc.). For this reason, I would expect that the method still should be prone to confounding and other causal spurious effects.
\end{point}
\begin{reply}
We acknowledge the reviewer’s concern regarding the theoretical basis of GCMC, particularly the adaptation from temporal to spatial domains and the introduction of the Intersectional Cardinality (IC) indicator. 

We would like to clarify that GCMC is grounded in generalized embedding theory \citep{Takens1981embedding,Mane1981embedding}, which posits that the dynamics of a system can be reconstructed from a single observable through delay-coordinate embedding. While traditional CCM employs time-delayed observations to construct state spaces, the GCMC method extends this principle to spatial domains by treating spatial lags as analogs of temporal lags. Although spatial lags differ in that multiple spatially adjacent values exist at a single lag step (as opposed to single temporal observations), this distinction is addressed through appropriate statistical aggregation before embedding. This spatial extension maintains the theoretical foundation established by dynamical systems theory and ensures that the shadow manifolds reconstructed from spatial neighbors can capture the intrinsic structure of the system.

Regarding the IC indicator, its introduction is motivated by the following principle: if two variables are driven by a shared underlying dynamical mechanism, their reconstructed state spaces (or manifolds) will exhibit structural overlap \citep{sugihara2012ccm,gao2023gccm}. This overlap allows for mutual cross mapping, as in CCM, but can also be directly reflected by the similarity in local neighborhood structures \citep{tao2023cmc}. GCMC leverages this idea by quantifying the degree of such structural overlap using the IC metric—i.e., the intersectional cardinality of nearest neighbors across reconstructed manifolds. This provides a compact, interpretable, and directional measure of dynamical causality, rooted in shared system dynamics.

To address the reviewer’s concern more explicitly, we have expanded the discussion of these theoretical underpinnings in Section~2.1 (Lines~87--91 and Lines~102--108) of the revised manuscript.
\end{reply}

\begin{point}
The authors designed evaluation experiments that in principle seems to be defensible (based on simulating spatial variables under different known causal model assumptions), and the results show that the method is able to distinguish some causal from non-causal effects. However, the evidence is very limited and preliminary: the benchmark comparison is only made with respect to their own previously developed method (GCCM) and with respect to non-causal methods like simple Pearson correlations and GD. That alone is not yet very convincing, because it leaves out all other approaches to causality detection that the method could be compared against.
\end{point}
\begin{reply}
We appreciate the reviewer’s careful examination and critical feedback. This comment helped us carefully reassess the logic and methodological rigor behind applying GCMC to real-world dataset.

In the original version of the manuscript, we underestimated the influence of the embedding dimension $E$ on the performance of GCMC. Specifically, we adopted a simple rule-of-thumb—setting $E = 2$ for smaller datasets and $E = 3$ for larger ones—without systematic justification. However, since the core of our method relies on the intersectional structure of nearest neighbors in reconstructed state spaces, using an inadequately low embedding dimension can lead to a high proportion of false nearest neighbors\citep{kennel1992fnn,ma2017cme}, thereby distorting the reconstructed manifold geometry and biasing the estimated causal strength.

To address this critical issue, we have now implemented the \textit{false nearest neighbors} method to determine appropriate embedding dimensions \citep{kennel1992fnn} for all variables prior to running GCMC or GCCM. The updated procedure is described in Lines~216--220 of the revised manuscript.

Furthermore, we have re-run all three case studies and the noise sensitivity analysis using the corrected embedding dimensions. The causal diagrams for the three case studies have been revised accordingly, incorporating both statistical significance testing and a stricter causal strength threshold to determine the presence of links. The updated results (presented in Figures~4–6 of the revised manuscript) no longer exhibit the inconsistencies criticized in the original submission. Specifically, in the first example, only causal links that are both statistically significant and have strength exceeding 0.1 are retained in the causal diagram (i.e., \texttt{crime}~$\rightarrow$~\texttt{hoval}, \texttt{inc}~$\rightarrow$~\texttt{hoval}; see Figure~\ref{fig_case1} below, or Figure~4 in the revised manuscript). In the second example, the causal link from temperature to elevation is excluded due to lack of statistical significance, while the link from population density to elevation is removed because its strength falls below the 0.10 threshold (see Figure~\ref{fig_case2} below or Figure~5 in the revised manuscript). In the third example, the link from NPP to temperature is excluded as well, owing to its lack of statistical significance (see Figure~\ref{fig_case3} below or Figure~6 in the revised manuscript). Lines~156–169 of the revised manuscript explain how causation is inferred from the causal strengths measured by GCMC.

\clearpage

\begin{figure}[htb!]
\centering
\includegraphics[width=1\textwidth]{./figure/figure4.png}
\caption{Causation among residential crime, house value, and household income (non-significant values are marked with $^\#$). \textbf{a-c} Spatial distributions of the three variables. \textbf{d} Correlations obtained by PCC, where high correlations indicate a moderately to strongly coupled scenario. \textbf{e} Stratified power computed by GD. \textbf{f} Causal strengths measured by GCCM. \textbf{g} Causal strengths measured by GCMC. \textbf{h} Conceptual diagram summarizing the inferred causal pathways based on GCMC estimates, where false positive links are excluded based on a causal strength threshold of $0.1$. Both GCMC and GCCM were configured with an embedding dimension \(E = 7\); the number of nearest neighbors used was \(k = 25\) for GCMC and \(k = 9\) for GCCM.}
\label{fig_case1}
\end{figure}

\begin{figure}[htb!]
\centering
\includegraphics[width=0.9\textwidth]{./figure/figure5.png}
\caption{Causation among elevation, temperature and population density (non-significant values are marked with $^\#$). \textbf{a} Correlations derived from PCC, which indicate that the dataset in this case represents a weak to moderate coupling state. \textbf{b} The stratified power quantified with GD. \textbf{c} Causal strengths measured by GCCM. \textbf{d} Causal strengths measured by GCMC. \textbf{e} Diagram of causal links constructed from GCMC results, with false positives filtered using a causal strength threshold of $0.1$. Both GCMC and GCCM were configured with an embedding dimension \(E = 10\); the number of nearest neighbors used was \(k = 200\) for GCMC and \(k = 12\) for GCCM.}
\label{fig_case2}
\end{figure}

\begin{figure}[htb!]
\centering
\includegraphics[width=0.9\textwidth]{./figure/figure6.png}
\caption{Causation among temperature, precipitation and net primary productivity (non-significant values are marked with $^\#$). \textbf{a} Correlations measured by PCC, which suggest that this case exhibits a high degree of coupling. \textbf{b} Stratified power computed by GD. \textbf{c} Causal strengths measured by GCCM. \textbf{d} Causal strengths measured by GCMC. \textbf{e} Causal links inferred by GCMC after filtering false positives using a causal strength threshold of $0.21$. Both GCMC and GCCM were configured with an embedding dimension \(E = 18\); the number of nearest neighbors used was \(k = 165\) for GCMC and \(k = 20\) for GCCM.}
\label{fig_case3}
\end{figure}

\clearpage

\end{reply}

\begin{point}
Furthermore, regarding the practicalities of the method itself, I see the problem that it makes heavy use of three modeling parameters that would need to be set in advance to make sure the method can distinguish causal from non-causal cases: a threshold parameter for effect strengths (in the experiments, this was 0.1, but how to determine this in general?), a significance threshold for the IC curve (which the authors don't seem to talk about), as well as a the spatial neighborhood lag parameter (k). All these parameters need to be set in a meaningful way (using what kind of knowledge?) before we can expect any meaningful outcomes. This shows that having a simple significance test for causality is illusory. In the end, you need a-priori knowledge to set the parameters right. This makes the whole approach a bit questionable I would say, though it might still be practically useful.
\end{point}
\begin{reply}
Thank you very much for your insightful comments. The main advantage of GCMC compared to GCCM lies in its relatively lower false positive rate and fewer causal omissions. For instance, in our own controlled benchmark tests where the causal structure is known, GCCM either misses causal paths or detects false positives in all three scenarios (as illustrated in Figure~3 of the revised manuscript, or as shown below). 

Regarding the specific case you mentioned in the NPP study, the issue you raised is resolved in GCMC when using the correct embedding dimension \( E \). In the revised results, the causal strength from NPP to temperature is statistically insignificant, leading us to exclude the causal link from NPP to temperature altogether.


\begin{figure}[ht!]
\centering
\includegraphics[width=1\textwidth]{./figure/figure3.png}
\caption{Benchmark evaluation of GCMC on canonical causal structures.
\textbf{(a)} Three simulated causal scenarios representing fundamental three-variable causal motifs. \textbf{(b)} Performance comparison of GCMC and GCCM on each scenario using estimated causal strengths (Values near edges). Edge types indicate causal inference correctness: true positives (black solid), false positives (red solid), true negatives (gray dashed), and false negatives (red dashed). Non-significant values are marked with $^\#$ and orange-highlighted thresholds indicate the optimal decision boundary.}
\label{fig_benchmark}
\end{figure}
\end{reply}

\begin{point}
Furthermore, the article lacks a clear research gap and a research question specific to this paper. Why precisely do we need another causal method? What is it that current causal (also spatially causal) methods cannot do, and why we need this method? What is the motivation for using this approach? This remains rather unclear, even after reading the article.
\end{point}
\begin{reply}

\end{reply}


\reviewersection 

\begin{point}
Title: This is a suitable, but dense title. Consider saying causal strengths of XXX, and maybe be more descriptive about what this cardinality is, or what settings this fits ( lattice data), to help readers identify the paper.
\end{point}
\begin{reply}
Thank you for raising this important question. By "causal principles", we refer primarily to well-established frameworks in causal inference, such as counterfactual reasoning \citep{rubin2005pom3}, structural causal models using directed acyclic graphs \citep{pearl2018scm4}, and predictive causality based on temporal or dynamical predictability \citep{granger1969grangertest,sugihara2012ccm}. These principles define causality as a relationship that can be identified under specific assumptions (e.g., ignorability, consistency) and not merely through statistical correlations. The GCMC method we introduce in this study is grounded in predictive causality, particularly as formalized in the empirical dynamic modeling (EDM) framework. Predictive causality—including Granger causality and its nonlinear extension such as convergent cross mapping (CCM)—posits that if variable $X$ causally influences $Y$, then the historical states of $Y$ should contain information about $X$. This form of causality is based on temporal predictability and dynamic information flow, which differs from traditional counterfactual or graph-based frameworks but remains a widely accepted causal concept in dynamic systems. To enhance clarity and avoid ambiguity, we have explicitly stated the specific causal principles referenced in this study in the revised manuscript (see Lines~37--42).
\end{reply}

\begin{point}
Intersectional cardinality may be a novel concept for the readership in IJGIS, and it would be worth toe explain it in more detail, including graphically. I suggest to plot of the AUC curves in Fig 2 should be used to discuss this more - explain what the curvature of the curve implies, etc. Similarly, the DeLong placement method is unfamiliar to GIScientists at large  I would argue, so it would be good to ease them in.
\end{point}
\begin{reply}
We acknowledge the reviewer’s comment regarding the need to define abbreviations upon first use. While ``AUC'' typically stands for \textit{area under the curve} in many analytical contexts, in our study it specifically refers to the \textit{area under the intersectional cardinality curve}. In GCMC, the difference between the AUCs under the null and alternative hypotheses of mutual cross mapping is further quantified as causal strength. To improve clarity, we have revised Line~70 of the manuscript to define AUC in full at its first appearance.
\end{reply}

\begin{point}
In Fig 2 we have the traditional attractors in a and b, typical for CCM. These are very hard to interpret, and I suspect the ones used here are ony mock ups. Maybe they should be labelled so.
\end{point}
\begin{reply}
We sincerely thank the reviewer for raising this thought-provoking and fundamental question. While both spatial association and causal association describe relationships between variables, they differ in nature and implication. Spatial association reflects the degree to which similar or dissimilar values co-occur in space, often measured by statistics such as Moran’s~$I$, Geary’s~$C$, or the GeoDetector $q$-statistic, and does not imply a directional or mechanistic relationship. In contrast, causal association implies that changes in one variable lead to or influence changes in another, underpinned by logical directionality and typically requiring additional assumptions or modeling frameworks to infer.

In the context of GCMC, although the estimation relies on patterns observed in reconstructed state spaces, it is designed to capture \textit{dynamical causality} by leveraging the intersectional cardinality of neighbors identified through mutual cross mapping. This approach provides directional insights by comparing the structure of mappings under the null and alternative hypotheses, combined with statistical significance testing and threshold filtering. These features distinguish GCMC from conventional measures of spatial association. To avoid misunderstanding, we have consistently used the term \textit{causal strength} throughout the manuscript to describe the results obtained from GCMC. Corresponding revisions have also been made to the manuscript title and related sections to emphasize this clarification.
\end{reply}

\begin{point}
Interpretation of the causal strength derived from the IC curve: I believe that these methods do not so much as quantify the treatment effect of a variable X_1, as much as enable to place the strength of this effect in context of other variables X_2..n.  In a traaditional interpretation, the derivation of a causal strength would be the multiplier by how much does an increase of X1 alter the outcome Y ( modulo confounders etc...). I am unsure whether the causala strength derived from the AUC of the Cardinality curve can be interpreted this way SO it may be a good mechanism for causal discovery, or even ranking of treatment effects of multiple variables, but I am not sure how it would relate to claims about absolute effect strength. Can you address the interpretation of the values you get from your GCMC measure's outcomes?
\end{point}
\begin{reply}
We thank the reviewer for highlighting this important point. Indeed, the way in which the explanatory variable is stratified can significantly affect the GD results. In all our case studies, we used the natural breaks method to discretize the explanatory variables into five classes when calculating the $q$-statistic (stratified power). This information has been explicitly reported in the revised manuscript (Lines~215--216).
\end{reply}

\baselineskip12pt
\bibliographystyle{elsarticle-harv} 
\bibliography{references.bib} 

\end{document}